\section{Introduction}
\label{sec:intro}

Isoparametric finite elements lose accuracy on distorted meshes because the
isoparametric map destroys the completeness of the approximation space in
Cartesian coordinates. Under angular distortion, serendipity elements no longer
contain all quadratic polynomials \citep{Lee1993effects,Rajendran2000completeness},
and under curved-edge distortion even full tensor-product (Lagrange) elements of
degree $p$ with a geometry map of degree $g$ are guaranteed to contain $\Pp_k(x,y)$ only
for $gk\le p$ \citep[cf.][]{Arnold2002Approximation}.
MacNeal's theorem \citep{MacNeal1992On} shows that, within the class of symmetric
(Bubnov--Galerkin) elements, distortion insensitivity and passing the patch test
cannot be combined in general. To circumvent this restriction,
\citet{Rajendran2003novel} proposed the unsymmetric finite element method (UFEM), a
Petrov--Galerkin formulation in which test and trial functions differ. Here, the test
functions are the usual parametric (isoparametric) functions, which guarantee
inter-element compatibility. The trial functions, on the other hand, are \emph{metric}
functions, i.e.\ polynomials in Cartesian coordinates that are interpolatory at the
element nodes. \Citet{Rajendran2010technique} formulated three sufficient conditions for
\emph{distortion immunity}: (i) compatible test functions, (ii) trial functions that are
complete in Cartesian coordinates and (iii) an element stiffness integrand free of
$\det\matr J$ in the denominator. Elements satisfying them reproduce every displacement
field of the completeness degree exactly on arbitrarily distorted meshes, including
meshes with $\det\matr J\le0$. The approach has been extended to frame-invariant
formulations in skew coordinates \citep{Yuan1993New,Ooi2008Remedies,Xie2016Modified}, to
analytical trial functions \citep{Cen2012shape,Cen2015unsymmetric,Cen2019Some}, to
orthotropic and incompatible-mode elements \citep{Shang2018airy,Huang2022Unsymmetric}, to
solids \citep{Ooi2004hexahedron}, to enhanced assumed strain formulations
\citep{Pfefferkorn2021mesh,Pfefferkorn2023hourglassing}, to geometrically nonlinear
problems \citep{Ooi2007extension,Li2018high} and to elastodynamics
\citep{Liew2006quadratic,Wang2022Temporal,Zaehringer2026On}, to name just a few. The
non-variational character of the formulation was analysed by
\citet{Prathap2006mesh,Prathap2007The}.

Almost all of this work concerns low-order elements (4 to 20 nodes). In the article at
hand, we consider the generalisation to spectral elements of arbitrary degree $p$. Here,
the parametric test functions are Lagrange polynomials on Gauss--Lobatto--Legendre (GLL)
nodes, and the metric trial functions are the unique functions of a
$(p+1)^2$-dimensional space that interpolate the GLL nodes in physical space. In the
standard construction, the metric space is $\Qp_p(s,t)$, the tensor-product polynomials in
the non-dimensional corotational skew coordinates $(s,t)$ of \citet{Xie2016Modified}, and
the shape functions are obtained by inverting a generalised Vandermonde matrix built from
the monomials $s^at^b$ \citep{Rajendran2003novel,Ooi2006unsymmetric}. At high order,
however, two difficulties of a different nature appear:
\begin{enumerate}
\item \emph{Ill-conditioning of the construction.} The monomial Vandermonde matrix
  is exponentially ill-conditioned in $p$ already on undistorted elements
  \citep{Gautschi1988Lower}. Consequently, the computed shape functions carry round-off
  errors of order $\cond(\matr V)\,u$, where $u$ denotes the unit round-off. These errors
  destroy, for instance, the symmetry of the stiffness matrix on affine elements and the
  Kronecker property on which the diagonal GLL mass matrix relies.
\item \emph{Ill-conditioning of the interpolation problem itself.} On non-affine
  elements, the images of the GLL nodes do not form a tensor grid in $(s,t)$. The
  Lebesgue constant of the interpolation in $\Qp_p(s,t)$ then grows exponentially with
  $p$ and with the amount of distortion. Note that this is a property of the space and of
  the node set, independent of how the shape functions are computed. The Lebesgue
  constant amplifies round-off, enters the interpolation error and -- as we show in
  \cref{sec:numerics} -- is accompanied by a loss of stability of the Petrov--Galerkin
  operator. That is to say, strongly complex and negative generalised eigenvalues appear,
  which make the semi-discrete equations of motion exponentially unstable, and the static
  solution diverges under mesh refinement.
\end{enumerate}

Therefore, the overarching goal of this paper is to present two improvements of the
metric shape functions that address these two difficulties separately.

The first improvement concerns the polynomial basis of the metric Vandermonde matrix
(\cref{sec:basis}). We show that the metric shape functions are independent of the
polynomial basis used to construct them (\cref{prop:invariance}). Thus, the basis is a
purely numerical choice, and we propose the Lagrange polynomials through the reference
GLL nodes in bounding-box scaled skew coordinates. On affine elements, the Vandermonde
matrix then becomes a permutation matrix ($\cond=1$ for every $p$). On distorted
elements, its condition number equals the Lebesgue function of the metric interpolation
on the reference grid of the bounding box, up to a factor bounded by the reference
Lebesgue constant (\cref{prop:kappa}). This traces the remaining ill-conditioning back to
the extrapolation of the interpolant beyond the element. Numerically, the condition
number at $p=16$ drops by three to four orders of magnitude relative to monomials. The
same analysis also shows that \emph{no} choice of basis can reduce the Lebesgue
constant.

The second improvement is a hybrid metric--parametric trial space (\cref{sec:hybrid}).
Distortion immunity requires only that the trial space contain the complete polynomials
$\Pp_p(x,y)$, i.e.\ $(p+1)(p+2)/2$ of the $(p+1)^2$ functions. Therefore, we propose the
trial space
\[
  \Tp_p = \Pp_p(x,y)\big|_{\Omega_e}\;\oplus\;\Cp_p\circ\vect F_e^{-1},
\]
in which only the complete polynomials are written in metric (skew) coordinates. The
remaining approximation, $\Cp_p$, is the $L^2$-orthogonal complement of
$\Pp_p(\xi,\eta)$ in $\Qp_p(\xi,\eta)$, written in the parametric coordinates, in
which the nodes form an exact tensor grid. We prove that $\Tp_p$ is complete,
affine invariant and interpolatory and that it coincides with the classical spectral
element space on every straight-edged quadrilateral. Furthermore, we derive an exact
representation of its interpolation operator as the GLL interpolant plus a correction
that is driven only by the completeness defect of the classical element
(\cref{prop:representation}). A factorisation of its Vandermonde determinant reduces
unisolvence to the non-singularity of a $\dim\Pp_p\times\dim\Pp_p$ matrix
(\cref{prop:factor}). Numerically, its Lebesgue constant stays at the level of GLL
interpolation for moderate curved-edge distortion and remains orders of magnitude below
that of $\Qp_p(s,t)$ for strong distortion. On all meshes tested, the eigenvalues with
negative real part of the $\Qp_p(s,t)$ element disappear, and the relative imaginary
parts drop from $\mathcal O(1)$ to below $2\cdot10^{-2}$.

The numerical studies (\cref{sec:numerics}) quantify conditioning, accuracy of the
shape functions, Lebesgue constants versus degree and distortion, patch tests on
distorted and folded meshes, $h$-convergence on meshes whose distortion persists
under refinement, eigenvalue spectra, transient stability and the thick-cylinder
benchmark of \citet{Rajendran2010technique}. In order to make the trade-off explicit,
they also document where the hybrid space is \emph{less} accurate than $\Qp_p(s,t)$
(static problems with smooth solutions at moderate degree and distortion). Note that only
linear plane elasticity on quadrilaterals is considered.

The outline of the remainder of the present article is as follows: First, the
unsymmetric spectral element formulation and the metric shape functions in skew
coordinates are briefly summarised in \cref{sec:formulation}. Thereafter, the polynomial
basis of the metric Vandermonde matrix is discussed in \cref{sec:basis}, before the
hybrid trial space is introduced and analysed in \cref{sec:hybrid}. The implementation is
outlined in \cref{sec:implementation}, and the numerical studies are presented in
\cref{sec:numerics}. Finally, the results are discussed in \cref{sec:discussion}, and the
most important conclusions are summarised in \cref{sec:conclusions}.
